\section{\texorpdfstring{$E_8$}{E8} from nine histories}
\label{sec:e8}

$E_8$ is the largest of the five exceptional simple Lie algebras: $248$ dimensions, $240$
roots, a symmetry so rigid that it appears, often unexpectedly, in string theory, in the
densest packing of spheres in eight dimensions, and in the physics of certain magnetic
materials. This section shows that one qutrit's nine operator histories build $E_8$ in a
way in which every piece has a meaning in the finite geometry --- and it records
carefully one famous coincidence that turned out to be a trap.

\begin{plainwords}
A Lie algebra is the ``infinitesimal'' version of a continuous symmetry: the set of all
directions in which you can nudge a system while keeping its rules intact, together with
the rule (the \emph{bracket}) for what happens when two nudges are combined in opposite
orders. Its \emph{roots} are the special directions that the symmetry's own ``clock''
(its maximal commuting part) rotates at definite rates. $E_8$ has $240$ of them.
\end{plainwords}

\subsection{Eighty plus eighty-four plus eighty-four}

Recall from Section~3 that the $80$ non-identity two-sided Weyl maps on one qutrit span
$\sll_9$, the Lie algebra of traceless $9\times9$ matrices acting on the nine operator
histories $\ket{i}\!\bra{j}$. Now add \emph{packets}: unordered triples of histories. There
are $\binom93=84$ of them, forming the representation $\Lambda^3\C^9$; their duals form
$\Lambda^6\C^9$, another $84$. A classical construction \cite{Adams1996,Baez2002} makes
\[
E_8 \;=\; \sll_9\;\oplus\;\Lambda^3\C^9\;\oplus\;\Lambda^6\C^9,\qquad
248=80+84+84 ,
\]
into a Lie algebra. On the nine history cells the brackets are literal packet operations:
$E_{ij}$ replaces cell $j$ by cell $i$ inside a packet; two disjoint triples bracket to the
complementary six-set; a triple and a six-set differing by one cell contract to an
$E_{ij}$; opposite packets return a Cartan direction. \tierC\
(\cert{w33\_20260924\_temporal\_su9\_e8\_compiler.py},
\cert{w33\_20260924\_temporal\_a8\_e8\_process\_bracket.py})

\subsection{Hesse lines are copies of \texorpdfstring{$A_2$}{A2}}

Place the nine histories on the affine plane $\F_3^2$ as in Figure~\ref{fig:hesse}. A
triple $S$ of cells labels the root
\[
w_S=\sum_{i\in S}e_i-\tfrac13\sum_{i=1}^9e_i ,\qquad w_S\cdot w_T=|S\cap T|-1 .
\]
(Indeed $|w_S|^2=3-2+1=2$.) The $84$ triples fall into two orbits of the affine group
$AGL(2,3)$ (order $432$): the $12$ \emph{lines} of the plane and the $72$ non-collinear
\emph{triangles}.

\begin{result}{Theorem 5.1 (four clocks, four $A_2$'s) \hfill \tierT\ \tierC}
The three lines of one parallel class give three roots with pairwise products $-1$ and
sum $0$; with their negatives they form a root system $A_2$. Lines of different classes
meet in exactly one cell, so their roots are orthogonal. The $12$ Hesse lines therefore
span a reflection-closed subsystem $4A_2\subset E_8$, one $A_2$ per clock, and the root
shell refines as
\[
240\;=\;\underbrace{72}_{A_8}\;+\;\underbrace{24}_{\text{Hesse }4A_2}\;+\;\underbrace{144}_{\text{triangles}} .
\]
\cert{w33\_20260924\_temporal\_hesse\_4a2\_a8\_bridge.py}
\end{result}

\begin{proof}
Disjoint triples have $|S\cap T|=0$, hence product $-1$; the three lines of a class
partition the nine cells, so $\sum w_S=\sum_i e_i-3\cdot\frac13\sum_ie_i=0$. Two
non-parallel lines of an affine plane meet in one point, giving product $1-1=0$.
\end{proof}

\begin{figure}[htbp]
\centering
\begin{tikzpicture}[scale=0.95]
  % four A2 hexagons, one per clock
  \foreach \c/\n/\x in {teal/1/0,sky/2/2.6,gold/3/5.2,sage/4/7.8}{
    \begin{scope}[xshift=\x cm]
      \foreach \k in {0,...,5}{
        \draw[\c,line width=0.9pt,-{Stealth[length=1.6mm]}] (0,0) -- (60*\k+30:0.9);}
      \draw[\c!35,line width=0.5pt] (30:0.9) \foreach \k in {1,...,5}{ -- (60*\k+30:0.9)} -- cycle;
      \fill[\c] (0,0) circle (1.3pt);
      \node[font=\sffamily\scriptsize\bfseries,text=\c!80!black] at (0,-1.25) {$A_2$, clock \n};
    \end{scope}}
  \node[font=\sffamily\scriptsize,text=slate] at (3.9,-1.85) {$4\times 6=24$ roots: the Hesse subsystem};
  % glue sectors
  \begin{scope}[xshift=11.3cm,yshift=0cm]
    \foreach \k in {0,...,7}{
      \pgfmathsetmacro{\a}{\k*45+22.5}
      \fill[rose!18,draw=rose!60] (\a:1.05) circle (0.33);
      \node[font=\tiny] at (\a:1.05) {$27$};}
    \node[font=\sffamily\scriptsize,align=center] at (0,0) {tetracode\\glue};
    \node[font=\sffamily\scriptsize,text=slate,align=center] at (0,-1.85) {$8$ nonzero codewords\\$\times\,27=216$ roots};
  \end{scope}
\end{tikzpicture}
\caption{$E_8$ assembled from four clocks. Each Hesse clock contributes one $A_2$ root
system (left). The remaining $216$ roots come in eight sectors of $27$, one for each
nonzero word of the tetracode (right): $240=4\cdot6+8\cdot27$.}
\label{fig:a24}
\end{figure}

\subsection{The glue is the tetracode, and the tetracode is the clock}

The lattice $E_8$ contains the sublattice $A_2^4$ with quotient $\F_3^2$, and the
well-known construction of $E_8$ from the tetracode \cite{ConwaySloane} reads it as
\[
E_8=\bigcup_{c\in\mathcal{T}}\bigl(c+A_2^{\,4}\bigr),\qquad
240=4\cdot6+8\cdot27 ,
\]
where $\mathcal{T}\subset\F_3^4$ is the tetracode: each of the eight nonzero codewords has
weight three, and contributes $3^3=27$ minimal vectors. Combined with Section~4, the
four glue coordinates are the four clocks and the eight nonzero sectors are the eight
oriented, nonzero linear clock readings. \tierT\ (classical) \ \tierC\ (the clock reading,
\cert{w33\_pass10946\_clock\_code\_cone\_objectwise.py})

\subsection{A trap: \texorpdfstring{$240=240$}{240 = 240}}
\label{sec:e8-trap}

The collinearity graph of $\W$ has $240$ edges, and $E_8$ has $240$ roots. For years the
repository looked for a natural bijection between them. The question has a precise and
negative answer.

\begin{result}{Theorem 5.2 (no equivariant bijection) \hfill \tierC}
The two groups of order $51\,840$ of Section~2 each do only half the job:
\begin{itemize}
\item $\PGSp(4,3)\cong W(E_6)$ acts faithfully and transitively on the $240$ edges of $\W$,
      but \emph{intransitively} on the $240$ roots of $E_8$;
\item $\Sp(4,3)$ acts faithfully and transitively on the $240$ roots, but its centre acts
      trivially on the edges, so its edge action is not faithful.
\end{itemize}
Hence no bijection between edges and roots is equivariant for a common group.
\cert{w33\_pass1020\_e8\_transitive\_51840.g}
\end{result}

What \emph{does} exist is better. The Weyl group of $E_8$ contains a regular element $w$
of order three. By Springer's theory \cite{Springer} its centraliser is the complex reflection group
$G_{32}\cong\Z_3\times\Sp(4,3)$ of order $155\,520$ --- the symmetry group of the
\emph{Witting polytope}, whose $240$ vertices in $\C^4$ are the $E_8$ roots.

\begin{result}{Theorem 5.3 (the $6:1$ fibration) \hfill \tierC}
The $240$ roots of $E_8$ fall into $40$ complex lines $\{\pm v,\pm\omega v,\pm\omega^2 v\}$,
each a copy of the six Eisenstein units. $\Sp(4,3)$ permutes the $40$ lines with
subdegrees $1,12,27$: the quotient is the collinearity graph of $\W$ --- on its
\emph{points}, not its lines. The fibre group is $\langle c^5\rangle\cong\Z_6$ for a
Coxeter element $c$ of $W(E_8)$.
\cert{w33\_pass1021\_e8\_fibration\_over\_forty.g}
\end{result}

\begin{figure}[htbp]
\centering
\begin{tikzpicture}[scale=1]
  % base: 40 points (sketch: the 1+12+27 view)
  \begin{scope}
    \fill[ink] (0,0) circle (2.2pt);
    \foreach \k in {0,...,11}{\fill[teal] ({30*\k}:0.9) circle (1.6pt);}
    \foreach \k in {0,...,26}{\fill[rose!70] ({360*\k/27+5}:1.6) circle (1.3pt);}
    \node[font=\sffamily\scriptsize,text=slate] at (0,-2.15) {base: the $40$ points of $\W$};
  \end{scope}
  % fibre over one point
  \draw[-{Stealth[length=2.5mm]},ink!60,line width=0.9pt] (3.1,1.3) to[bend right=25] node[above,font=\sffamily\scriptsize,pos=0.35]{project} (0.12,0.12);
  \begin{scope}[xshift=4.2cm,yshift=1.5cm]
    \foreach \k/\lab in {0/{v},1/{-\omega^2v},2/{\omega v},3/{-v},4/{\omega^2 v},5/{-\omega v}}{
      \fill[gold] (60*\k:1.0) circle (2.3pt);
      \node[font=\scriptsize] at (60*\k:1.42) {$\lab$};}
    \draw[gold!60] (0:1) \foreach \k in {1,...,5}{-- (60*\k:1)} -- cycle;
    \node[font=\sffamily\scriptsize,text=slate,align=center] at (0,-1.85) {fibre: $6$ roots $=$ one\\Witting diameter, $\Z_6$};
  \end{scope}
  \node[font=\sffamily\small,align=left,anchor=west] at (6.4,0.4)
   {$240\;=\;40\times6$\\[3pt]\scriptsize roots $\to$ points,\\\scriptsize $\Sp(4,3)$-equivariant};
\end{tikzpicture}
\caption{The correct relation between $E_8$ and $\W$: not edges to roots, but a
six-to-one map from roots to points, whose fibres are the Eisenstein unit hexagons.}
\label{fig:fibration}
\end{figure}

\begin{figure}[p]
\centering
\begin{tikzpicture}[scale=1.12]
  % generated by make_figures.py -- E8 Coxeter plane, 40 hexagonal fibres
\draw[ink!10,line width=0.3pt] (0,0) circle (0.753);
\draw[ink!10,line width=0.3pt] (0,0) circle (1.218);
\draw[ink!10,line width=0.3pt] (0,0) circle (1.497);
\draw[ink!10,line width=0.3pt] (0,0) circle (1.810);
\draw[ink!10,line width=0.3pt] (0,0) circle (2.225);
\draw[ink!10,line width=0.3pt] (0,0) circle (2.422);
\draw[ink!10,line width=0.3pt] (0,0) circle (2.929);
\draw[ink!10,line width=0.3pt] (0,0) circle (3.600);
\draw[teal!55,line width=0.45pt] (-1.721,-0.559) -- (-0.376,-1.770) -- (1.345,-1.211) -- (1.721,0.559) -- (0.376,1.770) -- (-1.345,1.211) -- cycle;
\fill[teal] (-1.721,-0.559) circle (1.25pt);
\fill[teal] (-0.376,-1.770) circle (1.25pt);
\fill[teal] (1.345,-1.211) circle (1.25pt);
\fill[teal] (1.721,0.559) circle (1.25pt);
\fill[teal] (0.376,1.770) circle (1.25pt);
\fill[teal] (-1.345,1.211) circle (1.25pt);
\draw[teal!55,line width=0.45pt] (-1.191,-0.253) -- (-0.376,-1.158) -- (0.815,-0.905) -- (1.191,0.253) -- (0.376,1.158) -- (-0.815,0.905) -- cycle;
\fill[teal] (-1.191,-0.253) circle (1.25pt);
\fill[teal] (-0.376,-1.158) circle (1.25pt);
\fill[teal] (0.815,-0.905) circle (1.25pt);
\fill[teal] (1.191,0.253) circle (1.25pt);
\fill[teal] (0.376,1.158) circle (1.25pt);
\fill[teal] (-0.815,0.905) circle (1.25pt);
\draw[sky!55,line width=0.45pt] (-1.218,-0.000) -- (-0.609,-1.055) -- (0.609,-1.055) -- (1.218,0.000) -- (0.609,1.055) -- (-0.609,1.055) -- cycle;
\fill[sky] (-1.218,-0.000) circle (1.25pt);
\fill[sky] (-0.609,-1.055) circle (1.25pt);
\fill[sky] (0.609,-1.055) circle (1.25pt);
\fill[sky] (1.218,0.000) circle (1.25pt);
\fill[sky] (0.609,1.055) circle (1.25pt);
\fill[sky] (-0.609,1.055) circle (1.25pt);
\draw[teal!55,line width=0.45pt] (-0.609,-0.442) -- (0.079,-0.748) -- (0.688,-0.306) -- (0.609,0.442) -- (-0.079,0.748) -- (-0.688,0.306) -- cycle;
\fill[teal] (-0.609,-0.442) circle (1.25pt);
\fill[teal] (0.079,-0.748) circle (1.25pt);
\fill[teal] (0.688,-0.306) circle (1.25pt);
\fill[teal] (0.609,0.442) circle (1.25pt);
\fill[teal] (-0.079,0.748) circle (1.25pt);
\fill[teal] (-0.688,0.306) circle (1.25pt);
\draw[sky!55,line width=0.45pt] (-0.688,-0.306) -- (-0.079,-0.748) -- (0.609,-0.442) -- (0.688,0.306) -- (0.079,0.748) -- (-0.609,0.442) -- cycle;
\fill[sky] (-0.688,-0.306) circle (1.25pt);
\fill[sky] (-0.079,-0.748) circle (1.25pt);
\fill[sky] (0.609,-0.442) circle (1.25pt);
\fill[sky] (0.688,0.306) circle (1.25pt);
\fill[sky] (0.079,0.748) circle (1.25pt);
\fill[sky] (-0.609,0.442) circle (1.25pt);
\draw[gold!55,line width=0.45pt] (-0.815,-0.905) -- (0.376,-1.158) -- (1.191,-0.253) -- (0.815,0.905) -- (-0.376,1.158) -- (-1.191,0.253) -- cycle;
\fill[gold] (-0.815,-0.905) circle (1.25pt);
\fill[gold] (0.376,-1.158) circle (1.25pt);
\fill[gold] (1.191,-0.253) circle (1.25pt);
\fill[gold] (0.815,0.905) circle (1.25pt);
\fill[gold] (-0.376,1.158) circle (1.25pt);
\fill[gold] (-1.191,0.253) circle (1.25pt);
\draw[sky!55,line width=0.45pt] (-1.345,-1.211) -- (0.376,-1.770) -- (1.721,-0.559) -- (1.345,1.211) -- (-0.376,1.770) -- (-1.721,0.559) -- cycle;
\fill[sky] (-1.345,-1.211) circle (1.25pt);
\fill[sky] (0.376,-1.770) circle (1.25pt);
\fill[sky] (1.721,-0.559) circle (1.25pt);
\fill[sky] (1.345,1.211) circle (1.25pt);
\fill[sky] (-0.376,1.770) circle (1.25pt);
\fill[sky] (-1.721,0.559) circle (1.25pt);
\draw[gold!55,line width=0.45pt] (-1.567,-0.905) -- (0.000,-1.810) -- (1.567,-0.905) -- (1.567,0.905) -- (-0.000,1.810) -- (-1.567,0.905) -- cycle;
\fill[gold] (-1.567,-0.905) circle (1.25pt);
\fill[gold] (0.000,-1.810) circle (1.25pt);
\fill[gold] (1.567,-0.905) circle (1.25pt);
\fill[gold] (1.567,0.905) circle (1.25pt);
\fill[gold] (-0.000,1.810) circle (1.25pt);
\fill[gold] (-1.567,0.905) circle (1.25pt);
\draw[teal!55,line width=0.45pt] (-2.098,-1.211) -- (0.000,-2.422) -- (2.098,-1.211) -- (2.098,1.211) -- (-0.000,2.422) -- (-2.098,1.211) -- cycle;
\fill[teal] (-2.098,-1.211) circle (1.25pt);
\fill[teal] (0.000,-2.422) circle (1.25pt);
\fill[teal] (2.098,-1.211) circle (1.25pt);
\fill[teal] (2.098,1.211) circle (1.25pt);
\fill[teal] (-0.000,2.422) circle (1.25pt);
\fill[teal] (-2.098,1.211) circle (1.25pt);
\draw[teal!55,line width=0.45pt] (-2.912,-0.306) -- (-1.191,-2.675) -- (1.721,-2.369) -- (2.912,0.306) -- (1.191,2.675) -- (-1.721,2.369) -- cycle;
\fill[teal] (-2.912,-0.306) circle (1.25pt);
\fill[teal] (-1.191,-2.675) circle (1.25pt);
\fill[teal] (1.721,-2.369) circle (1.25pt);
\fill[teal] (2.912,0.306) circle (1.25pt);
\fill[teal] (1.191,2.675) circle (1.25pt);
\fill[teal] (-1.721,2.369) circle (1.25pt);
\draw[sky!55,line width=0.45pt] (-1.721,-2.369) -- (1.191,-2.675) -- (2.912,-0.306) -- (1.721,2.369) -- (-1.191,2.675) -- (-2.912,0.306) -- cycle;
\fill[sky] (-1.721,-2.369) circle (1.25pt);
\fill[sky] (1.191,-2.675) circle (1.25pt);
\fill[sky] (2.912,-0.306) circle (1.25pt);
\fill[sky] (1.721,2.369) circle (1.25pt);
\fill[sky] (-1.191,2.675) circle (1.25pt);
\fill[sky] (-2.912,0.306) circle (1.25pt);
\draw[sky!55,line width=0.45pt] (-2.409,-0.253) -- (-0.985,-2.213) -- (1.424,-1.960) -- (2.409,0.253) -- (0.985,2.213) -- (-1.424,1.960) -- cycle;
\fill[sky] (-2.409,-0.253) circle (1.25pt);
\fill[sky] (-0.985,-2.213) circle (1.25pt);
\fill[sky] (1.424,-1.960) circle (1.25pt);
\fill[sky] (2.409,0.253) circle (1.25pt);
\fill[sky] (0.985,2.213) circle (1.25pt);
\fill[sky] (-1.424,1.960) circle (1.25pt);
\draw[gold!55,line width=0.45pt] (-0.504,-0.559) -- (0.233,-0.716) -- (0.736,-0.156) -- (0.504,0.559) -- (-0.233,0.716) -- (-0.736,0.156) -- cycle;
\fill[gold] (-0.504,-0.559) circle (1.25pt);
\fill[gold] (0.233,-0.716) circle (1.25pt);
\fill[gold] (0.736,-0.156) circle (1.25pt);
\fill[gold] (0.504,0.559) circle (1.25pt);
\fill[gold] (-0.233,0.716) circle (1.25pt);
\fill[gold] (-0.736,0.156) circle (1.25pt);
\draw[sage!55,line width=0.45pt] (-1.064,-1.464) -- (0.736,-1.653) -- (1.800,-0.189) -- (1.064,1.464) -- (-0.736,1.653) -- (-1.800,0.189) -- cycle;
\fill[sage] (-1.064,-1.464) circle (1.25pt);
\fill[sage] (0.736,-1.653) circle (1.25pt);
\fill[sage] (1.800,-0.189) circle (1.25pt);
\fill[sage] (1.064,1.464) circle (1.25pt);
\fill[sage] (-0.736,1.653) circle (1.25pt);
\fill[sage] (-1.800,0.189) circle (1.25pt);
\draw[teal!55,line width=0.45pt] (-1.112,-1.002) -- (0.311,-1.464) -- (1.424,-0.463) -- (1.112,1.002) -- (-0.311,1.464) -- (-1.424,0.463) -- cycle;
\fill[teal] (-1.112,-1.002) circle (1.25pt);
\fill[teal] (0.311,-1.464) circle (1.25pt);
\fill[teal] (1.424,-0.463) circle (1.25pt);
\fill[teal] (1.112,1.002) circle (1.25pt);
\fill[teal] (-0.311,1.464) circle (1.25pt);
\fill[teal] (-1.424,0.463) circle (1.25pt);
\draw[sage!55,line width=0.45pt] (-0.985,-0.716) -- (0.127,-1.211) -- (1.112,-0.495) -- (0.985,0.716) -- (-0.127,1.211) -- (-1.112,0.495) -- cycle;
\fill[sage] (-0.985,-0.716) circle (1.25pt);
\fill[sage] (0.127,-1.211) circle (1.25pt);
\fill[sage] (1.112,-0.495) circle (1.25pt);
\fill[sage] (0.985,0.716) circle (1.25pt);
\fill[sage] (-0.127,1.211) circle (1.25pt);
\fill[sage] (-1.112,0.495) circle (1.25pt);
\draw[teal!55,line width=0.45pt] (-1.800,-1.308) -- (0.233,-2.213) -- (2.033,-0.905) -- (1.800,1.308) -- (-0.233,2.213) -- (-2.033,0.905) -- cycle;
\fill[teal] (-1.800,-1.308) circle (1.25pt);
\fill[teal] (0.233,-2.213) circle (1.25pt);
\fill[teal] (2.033,-0.905) circle (1.25pt);
\fill[teal] (1.800,1.308) circle (1.25pt);
\fill[teal] (-0.233,2.213) circle (1.25pt);
\fill[teal] (-2.033,0.905) circle (1.25pt);
\draw[sage!55,line width=0.45pt] (-0.753,-0.000) -- (-0.376,-0.652) -- (0.376,-0.652) -- (0.753,0.000) -- (0.376,0.652) -- (-0.376,0.652) -- cycle;
\fill[sage] (-0.753,-0.000) circle (1.25pt);
\fill[sage] (-0.376,-0.652) circle (1.25pt);
\fill[sage] (0.376,-0.652) circle (1.25pt);
\fill[sage] (0.753,0.000) circle (1.25pt);
\fill[sage] (0.376,0.652) circle (1.25pt);
\fill[sage] (-0.376,0.652) circle (1.25pt);
\draw[gold!55,line width=0.45pt] (-2.536,-1.464) -- (0.000,-2.929) -- (2.536,-1.464) -- (2.536,1.464) -- (-0.000,2.929) -- (-2.536,1.464) -- cycle;
\fill[gold] (-2.536,-1.464) circle (1.25pt);
\fill[gold] (0.000,-2.929) circle (1.25pt);
\fill[gold] (2.536,-1.464) circle (1.25pt);
\fill[gold] (2.536,1.464) circle (1.25pt);
\fill[gold] (-0.000,2.929) circle (1.25pt);
\fill[gold] (-2.536,1.464) circle (1.25pt);
\draw[teal!55,line width=0.45pt] (-3.289,-1.464) -- (-0.376,-3.580) -- (2.912,-2.116) -- (3.289,1.464) -- (0.376,3.580) -- (-2.912,2.116) -- cycle;
\fill[teal] (-3.289,-1.464) circle (1.25pt);
\fill[teal] (-0.376,-3.580) circle (1.25pt);
\fill[teal] (2.912,-2.116) circle (1.25pt);
\fill[teal] (3.289,1.464) circle (1.25pt);
\fill[teal] (0.376,3.580) circle (1.25pt);
\fill[teal] (-2.912,2.116) circle (1.25pt);
\draw[sky!55,line width=0.45pt] (-1.296,-0.748) -- (0.000,-1.497) -- (1.296,-0.748) -- (1.296,0.748) -- (-0.000,1.497) -- (-1.296,0.748) -- cycle;
\fill[sky] (-1.296,-0.748) circle (1.25pt);
\fill[sky] (0.000,-1.497) circle (1.25pt);
\fill[sky] (1.296,-0.748) circle (1.25pt);
\fill[sky] (1.296,0.748) circle (1.25pt);
\fill[sky] (-0.000,1.497) circle (1.25pt);
\fill[sky] (-1.296,0.748) circle (1.25pt);
\draw[rose!55,line width=0.45pt] (-1.112,-0.495) -- (-0.127,-1.211) -- (0.985,-0.716) -- (1.112,0.495) -- (0.127,1.211) -- (-0.985,0.716) -- cycle;
\fill[rose] (-1.112,-0.495) circle (1.25pt);
\fill[rose] (-0.127,-1.211) circle (1.25pt);
\fill[rose] (0.985,-0.716) circle (1.25pt);
\fill[rose] (1.112,0.495) circle (1.25pt);
\fill[rose] (0.127,1.211) circle (1.25pt);
\fill[rose] (-0.985,0.716) circle (1.25pt);
\draw[sky!55,line width=0.45pt] (-2.033,-0.905) -- (-0.233,-2.213) -- (1.800,-1.308) -- (2.033,0.905) -- (0.233,2.213) -- (-1.800,1.308) -- cycle;
\fill[sky] (-2.033,-0.905) circle (1.25pt);
\fill[sky] (-0.233,-2.213) circle (1.25pt);
\fill[sky] (1.800,-1.308) circle (1.25pt);
\fill[sky] (2.033,0.905) circle (1.25pt);
\fill[sky] (0.233,2.213) circle (1.25pt);
\fill[sky] (-1.800,1.308) circle (1.25pt);
\draw[gold!55,line width=0.45pt] (-0.880,-1.211) -- (0.609,-1.368) -- (1.489,-0.156) -- (0.880,1.211) -- (-0.609,1.368) -- (-1.489,0.156) -- cycle;
\fill[gold] (-0.880,-1.211) circle (1.25pt);
\fill[gold] (0.609,-1.368) circle (1.25pt);
\fill[gold] (1.489,-0.156) circle (1.25pt);
\fill[gold] (0.880,1.211) circle (1.25pt);
\fill[gold] (-0.609,1.368) circle (1.25pt);
\fill[gold] (-1.489,0.156) circle (1.25pt);
\draw[sage!55,line width=0.45pt] (-2.785,-0.905) -- (-0.609,-2.865) -- (2.176,-1.960) -- (2.785,0.905) -- (0.609,2.865) -- (-2.176,1.960) -- cycle;
\fill[sage] (-2.785,-0.905) circle (1.25pt);
\fill[sage] (-0.609,-2.865) circle (1.25pt);
\fill[sage] (2.176,-1.960) circle (1.25pt);
\fill[sage] (2.785,0.905) circle (1.25pt);
\fill[sage] (0.609,2.865) circle (1.25pt);
\fill[sage] (-2.176,1.960) circle (1.25pt);
\draw[sky!55,line width=0.45pt] (-3.521,-0.748) -- (-1.112,-3.424) -- (2.409,-2.675) -- (3.521,0.748) -- (1.112,3.424) -- (-2.409,2.675) -- cycle;
\fill[sky] (-3.521,-0.748) circle (1.25pt);
\fill[sky] (-1.112,-3.424) circle (1.25pt);
\fill[sky] (2.409,-2.675) circle (1.25pt);
\fill[sky] (3.521,0.748) circle (1.25pt);
\fill[sky] (1.112,3.424) circle (1.25pt);
\fill[sky] (-2.409,2.675) circle (1.25pt);
\draw[gold!55,line width=0.45pt] (-1.800,-1.621) -- (0.504,-2.369) -- (2.304,-0.748) -- (1.800,1.621) -- (-0.504,2.369) -- (-2.304,0.748) -- cycle;
\fill[gold] (-1.800,-1.621) circle (1.25pt);
\fill[gold] (0.504,-2.369) circle (1.25pt);
\fill[gold] (2.304,-0.748) circle (1.25pt);
\fill[gold] (1.800,1.621) circle (1.25pt);
\fill[gold] (-0.504,2.369) circle (1.25pt);
\fill[gold] (-2.304,0.748) circle (1.25pt);
\draw[rose!55,line width=0.45pt] (-1.800,-0.189) -- (-0.736,-1.653) -- (1.064,-1.464) -- (1.800,0.189) -- (0.736,1.653) -- (-1.064,1.464) -- cycle;
\fill[rose] (-1.800,-0.189) circle (1.25pt);
\fill[rose] (-0.736,-1.653) circle (1.25pt);
\fill[rose] (1.064,-1.464) circle (1.25pt);
\fill[rose] (1.800,0.189) circle (1.25pt);
\fill[rose] (0.736,1.653) circle (1.25pt);
\fill[rose] (-1.064,1.464) circle (1.25pt);
\draw[sage!55,line width=0.45pt] (-1.424,-0.463) -- (-0.311,-1.464) -- (1.112,-1.002) -- (1.424,0.463) -- (0.311,1.464) -- (-1.112,1.002) -- cycle;
\fill[sage] (-1.424,-0.463) circle (1.25pt);
\fill[sage] (-0.311,-1.464) circle (1.25pt);
\fill[sage] (1.112,-1.002) circle (1.25pt);
\fill[sage] (1.424,0.463) circle (1.25pt);
\fill[sage] (0.311,1.464) circle (1.25pt);
\fill[sage] (-1.112,1.002) circle (1.25pt);
\draw[gold!55,line width=0.45pt] (-1.489,-1.653) -- (0.688,-2.116) -- (2.176,-0.463) -- (1.489,1.653) -- (-0.688,2.116) -- (-2.176,0.463) -- cycle;
\fill[gold] (-1.489,-1.653) circle (1.25pt);
\fill[gold] (0.688,-2.116) circle (1.25pt);
\fill[gold] (2.176,-0.463) circle (1.25pt);
\fill[gold] (1.489,1.653) circle (1.25pt);
\fill[gold] (-0.688,2.116) circle (1.25pt);
\fill[gold] (-2.176,0.463) circle (1.25pt);
\draw[sage!55,line width=0.45pt] (-2.176,-0.463) -- (-0.688,-2.116) -- (1.489,-1.653) -- (2.176,0.463) -- (0.688,2.116) -- (-1.489,1.653) -- cycle;
\fill[sage] (-2.176,-0.463) circle (1.25pt);
\fill[sage] (-0.688,-2.116) circle (1.25pt);
\fill[sage] (1.489,-1.653) circle (1.25pt);
\fill[sage] (2.176,0.463) circle (1.25pt);
\fill[sage] (0.688,2.116) circle (1.25pt);
\fill[sage] (-1.489,1.653) circle (1.25pt);
\draw[gold!55,line width=0.45pt] (-3.600,-0.000) -- (-1.800,-3.118) -- (1.800,-3.118) -- (3.600,0.000) -- (1.800,3.118) -- (-1.800,3.118) -- cycle;
\fill[gold] (-3.600,-0.000) circle (1.25pt);
\fill[gold] (-1.800,-3.118) circle (1.25pt);
\fill[gold] (1.800,-3.118) circle (1.25pt);
\fill[gold] (3.600,0.000) circle (1.25pt);
\fill[gold] (1.800,3.118) circle (1.25pt);
\fill[gold] (-1.800,3.118) circle (1.25pt);
\draw[rose!55,line width=0.45pt] (-2.225,-0.000) -- (-1.112,-1.927) -- (1.112,-1.927) -- (2.225,0.000) -- (1.112,1.927) -- (-1.112,1.927) -- cycle;
\fill[rose] (-2.225,-0.000) circle (1.25pt);
\fill[rose] (-1.112,-1.927) circle (1.25pt);
\fill[rose] (1.112,-1.927) circle (1.25pt);
\fill[rose] (2.225,0.000) circle (1.25pt);
\fill[rose] (1.112,1.927) circle (1.25pt);
\fill[rose] (-1.112,1.927) circle (1.25pt);
\draw[sage!55,line width=0.45pt] (-1.424,-1.960) -- (0.985,-2.213) -- (2.409,-0.253) -- (1.424,1.960) -- (-0.985,2.213) -- (-2.409,0.253) -- cycle;
\fill[sage] (-1.424,-1.960) circle (1.25pt);
\fill[sage] (0.985,-2.213) circle (1.25pt);
\fill[sage] (2.409,-0.253) circle (1.25pt);
\fill[sage] (1.424,1.960) circle (1.25pt);
\fill[sage] (-0.985,2.213) circle (1.25pt);
\fill[sage] (-2.409,0.253) circle (1.25pt);
\draw[rose!55,line width=0.45pt] (-0.736,-0.156) -- (-0.233,-0.716) -- (0.504,-0.559) -- (0.736,0.156) -- (0.233,0.716) -- (-0.504,0.559) -- cycle;
\fill[rose] (-0.736,-0.156) circle (1.25pt);
\fill[rose] (-0.233,-0.716) circle (1.25pt);
\fill[rose] (0.504,-0.559) circle (1.25pt);
\fill[rose] (0.736,0.156) circle (1.25pt);
\fill[rose] (0.233,0.716) circle (1.25pt);
\fill[rose] (-0.504,0.559) circle (1.25pt);
\draw[rose!55,line width=0.45pt] (-2.176,-1.960) -- (0.609,-2.865) -- (2.785,-0.905) -- (2.176,1.960) -- (-0.609,2.865) -- (-2.785,0.905) -- cycle;
\fill[rose] (-2.176,-1.960) circle (1.25pt);
\fill[rose] (0.609,-2.865) circle (1.25pt);
\fill[rose] (2.785,-0.905) circle (1.25pt);
\fill[rose] (2.176,1.960) circle (1.25pt);
\fill[rose] (-0.609,2.865) circle (1.25pt);
\fill[rose] (-2.785,0.905) circle (1.25pt);
\draw[rose!55,line width=0.45pt] (-1.489,-0.156) -- (-0.609,-1.368) -- (0.880,-1.211) -- (1.489,0.156) -- (0.609,1.368) -- (-0.880,1.211) -- cycle;
\fill[rose] (-1.489,-0.156) circle (1.25pt);
\fill[rose] (-0.609,-1.368) circle (1.25pt);
\fill[rose] (0.880,-1.211) circle (1.25pt);
\fill[rose] (1.489,0.156) circle (1.25pt);
\fill[rose] (0.609,1.368) circle (1.25pt);
\fill[rose] (-0.880,1.211) circle (1.25pt);
\draw[sage!55,line width=0.45pt] (-2.409,-2.675) -- (1.112,-3.424) -- (3.521,-0.748) -- (2.409,2.675) -- (-1.112,3.424) -- (-3.521,0.748) -- cycle;
\fill[sage] (-2.409,-2.675) circle (1.25pt);
\fill[sage] (1.112,-3.424) circle (1.25pt);
\fill[sage] (3.521,-0.748) circle (1.25pt);
\fill[sage] (2.409,2.675) circle (1.25pt);
\fill[sage] (-1.112,3.424) circle (1.25pt);
\fill[sage] (-3.521,0.748) circle (1.25pt);
\draw[rose!55,line width=0.45pt] (-2.304,-0.748) -- (-0.504,-2.369) -- (1.800,-1.621) -- (2.304,0.748) -- (0.504,2.369) -- (-1.800,1.621) -- cycle;
\fill[rose] (-2.304,-0.748) circle (1.25pt);
\fill[rose] (-0.504,-2.369) circle (1.25pt);
\fill[rose] (1.800,-1.621) circle (1.25pt);
\fill[rose] (2.304,0.748) circle (1.25pt);
\fill[rose] (0.504,2.369) circle (1.25pt);
\fill[rose] (-1.800,1.621) circle (1.25pt);
\draw[rose!55,line width=0.45pt] (-2.912,-2.116) -- (0.376,-3.580) -- (3.289,-1.464) -- (2.912,2.116) -- (-0.376,3.580) -- (-3.289,1.464) -- cycle;
\fill[rose] (-2.912,-2.116) circle (1.25pt);
\fill[rose] (0.376,-3.580) circle (1.25pt);
\fill[rose] (3.289,-1.464) circle (1.25pt);
\fill[rose] (2.912,2.116) circle (1.25pt);
\fill[rose] (-0.376,3.580) circle (1.25pt);
\fill[rose] (-3.289,1.464) circle (1.25pt);

  \fill[ink] (0,0) circle (1.2pt);
\end{tikzpicture}
\caption{The $240$ roots of $E_8$ projected to the Coxeter plane of a Coxeter element $c$
(order $30$), which acts on the plane as a rotation by $12^\circ$. The roots lie on eight
rings of thirty. The orbits of $c^5$ --- rotation by $60^\circ$ --- are regular hexagons,
five per ring: these $40$ hexagons are the $40$ fibres of Theorem~5.3, one for each point of
$\W$. Each hexagon is $\{\pm v,\pm\omega v,\pm\omega^2v\}$, since $c^{15}=-1$ and $c^{10}$ has
order three. All coordinates computed (\texttt{make\_figures.py}).}
\label{fig:coxeter}
\end{figure}

\begin{physical}[What it means physically: a photon's four-dimensional state space]
The $40$ Witting diameters are $40$ rays in $\C^4$ --- the state space of a photon's path
and polarisation, or of two qubits. Two rays are orthogonal exactly when the corresponding
points of $\W$ are collinear. So the same geometry is drawn twice: by \emph{states} of a
four-level system (the Witting rays) and by \emph{operators} on two qutrits (the Pauli
classes). The photonic Holonet papers call this operator--operand duality.
\tierC\ (\cert{bt817\_self\_entangled\_photon\_atlas.py}, \cert{bt821\_operator\_operand\_duality.py})
\end{physical}

\subsection{\texorpdfstring{$E_6\times A_2$}{E6 x A2}, the twenty-seven, and a cubic}

Under its subalgebra $E_6\oplus A_2$,
\[
E_8=(\mathbf{78},\mathbf1)\oplus(\mathbf1,\mathbf8)\oplus(\mathbf{27},\mathbf3)\oplus(\overline{\mathbf{27}},\overline{\mathbf3}),
\]
which is a $\Z_3$-grading $248=86+81+81$. The $27$ of $E_6$ carries a unique invariant
cubic form $N$; in the repository's committed basis it has exactly $45$ signed terms
(\emph{triads}), each coordinate in $5$ triads and each pair in at most one ---
precisely the incidence of the $27$ lines and $45$ tritangent planes of a smooth cubic
surface, whose symmetry group is $W(E_6)$, the automorphism group of $\W$ \cite{Garibaldi}.
\tierT\ (classical) \ \tierC\ (\cert{canonical\_su3\_gauge\_and\_cubic.json})

\begin{physical}
In grand unified theories built on $E_6$, the $\mathbf{27}$ holds one generation of
matter: under $SO(10)$ it splits $\mathbf{27}=\mathbf{16}+\mathbf{10}+\mathbf1$ ---
sixteen fermions (including a right-handed neutrino), a ten of Higgs-type fields, and a
singlet. The cubic $N$ is then the unique renormalisable Yukawa coupling
$\mathbf{27}^3$. Section~8 shows this same split arising as scalar, vector and spinor of a
ten-dimensional Lorentz group.
\end{physical}

\subsection{Clock gradings of \texorpdfstring{$E_8$}{E8}}

A clock gives more than a $\Z_3$-grading: once an origin and an orientation are chosen it
gives an \emph{integer} grading. On the committed Chevalley basis:

\begin{result}{Theorem 5.4 (Hesse clocks grade $E_8$) \hfill \tierC}
\begin{enumerate}[label=(\roman*),leftmargin=1.5em]
\item A single Hesse line (a ``tick'') induces the \emph{contact grading}
      \[ E_8=\g_{-2}\oplus\g_{-1}\oplus\g_0\oplus\g_1\oplus\g_2,\qquad 1+56+134+56+1, \]
      with $\g_0=E_7\oplus\C$.
\item An oriented striation (three parallel ticks) induces the \emph{$|3|$-grading}
      \[ 2+27+54+82+54+27+2, \]
      with Levi $\g_0=E_6\oplus A_1\oplus\C$, $\g_{\pm1}=\mathbf{27}\otimes\mathbf2$,
      $\g_{\pm2}=\mathbf{27}$, $\g_{\pm3}=\mathbf2$.
\item The repository's frozen $\Z_3$-grading $86+81+81$ is exactly this integer grading
      reduced modulo three; all $8\,347$ nonzero committed structure constants obey the
      integer degree law.
\end{enumerate}
\cert{w33\_20260925\_hesse\_clock\_e8\_gradings.py},
\cert{w33\_20260925\_e8\_parabolic\_cubic\_clock\_lift.py}
\end{result}

\begin{figure}[htbp]
\centering
\begin{tikzpicture}[yscale=0.028,xscale=1]
  % contact grading
  \begin{scope}
    \foreach \d/\h/\c in {-2/1/ink,-1/56/teal,0/134/gold,1/56/teal,2/1/ink}{
      \fill[\c!70] (\d*0.9-0.32,0) rectangle (\d*0.9+0.32,\h);
      \node[font=\scriptsize,anchor=south] at (\d*0.9,\h+2) {$\h$};
      \node[font=\scriptsize,anchor=north] at (\d*0.9,-3) {$\d$};}
    \node[font=\sffamily\scriptsize,text=slate] at (0,-22) {contact grading (one tick)};
    \node[font=\scriptsize,text=white] at (0,40) {$E_7$};
    \node[font=\scriptsize,text=white] at (0,28) {$+\C$};
    \node[font=\tiny,text=white,rotate=90] at (-0.9,26) {Freudenthal};
    \node[font=\tiny,text=white,rotate=90] at (0.9,26) {Freudenthal};
  \end{scope}
  % |3| grading
  \begin{scope}[xshift=7.2cm]
    \foreach \d/\h/\c in {-3/2/ink,-2/27/rose,-1/54/teal,0/82/gold,1/54/teal,2/27/rose,3/2/ink}{
      \fill[\c!70] (\d*0.9-0.32,0) rectangle (\d*0.9+0.32,\h);
      \node[font=\scriptsize,anchor=south] at (\d*0.9,\h+2) {$\h$};
      \node[font=\scriptsize,anchor=north] at (\d*0.9,-3) {$\d$};}
    \node[font=\sffamily\scriptsize,text=slate] at (0,-22) {$|3|$-grading (oriented striation)};
    \node[font=\tiny,text=white,align=center] at (0,33) {$E_6$\\$A_1$\\$\C$};
  \end{scope}
\end{tikzpicture}
\caption{Two ways a clock slices $E_8$. Heights are dimensions. On the left, the
Freudenthal ``$56$'' of Section~6 appears at degree $\pm1$; on the right, the $27$ of
$E_6$ appears doubled at degree $\pm1$ and single at degree $\pm2$.}
\label{fig:gradings}
\end{figure}

This parabolic spectrum was found independently by Kraft, Regeta and Zimmermann in a
study of small varieties with group actions \cite{KRZ2020}. What is repository-specific is
the identification with the Hesse clocks and the committed structure constants.

\subsection{The cubic ladder: three steps to the top}

In the $|3|$-grading, positive degrees form a nilpotent algebra
$\g_+=\g_1\oplus\g_2\oplus\g_3$ of dimension $54+27+2=83$. It is generated by $\g_1$, and
the way it fills up is governed by the $E_6$ cubic.

\begin{result}{Theorem 5.5 (the bracket ladder) \hfill \tierC}
On the committed Chevalley bracket:
\begin{enumerate}[label=(\roman*),leftmargin=1.5em]
\item $[\g_1,\g_1]=\g_2$: $270$ nonzero pairs hit all $27$ grade-two roots, each ten times;
\item $[\g_1,\g_2]=\g_3$: $54$ nonzero pairs hit both top roots, each $27$ times;
\item the nonzero nested triple brackets exhibit exactly the $45$ triads of the cubic,
      each $12$ times, and for every signed triad $d_{ijk}$,
      $[[e_{(i,0)},e_{(j,1)}],e_{(k,a)}]=-d_{ijk}\,T_a$ for $a\in\{0,1\}$.
\end{enumerate}
So the positive part grows $54\to81\to83$, and the committed cubic is exactly the rule by
which three first-order moves reach the two-dimensional top.
\cert{w33\_20260925\_cubic\_clock\_bracket\_ladder.py}
\end{result}

\begin{figure}[htbp]
\centering
\begin{tikzpicture}[>={Stealth[length=2.3mm]},
   blk/.style={draw=#1,fill=#1!12,rounded corners=3pt,align=center,font=\sffamily\scriptsize,minimum height=9mm}]
  \node[blk=teal,minimum width=36mm] (g1) at (0,0) {$\g_1=\mathbf{27}\otimes\mathbf2$\\$54$ first-order moves};
  \node[blk=rose,minimum width=26mm] (g2) at (4.6,0) {$\g_2=\mathbf{27}$\\$27$ histories};
  \node[blk=ink,minimum width=18mm] (g3) at (8.3,0) {$\g_3=\mathbf2$\\top doublet};
  \draw[->,line width=1pt] (g1) -- node[above,font=\scriptsize]{$[\g_1,\g_1]$} node[below,font=\tiny,text=slate]{$270$ pairs, $\times10$} (g2);
  \draw[->,line width=1pt] (g2) -- node[above,font=\scriptsize]{$[\g_1,\cdot\,]$} node[below,font=\tiny,text=slate]{$54$ pairs, $\times27$} (g3);
  \node[font=\sffamily\scriptsize,text=slate] at (0,-1.1) {cumulative $54$};
  \node[font=\sffamily\scriptsize,text=slate] at (4.6,-1.1) {cumulative $81$};
  \node[font=\sffamily\scriptsize,text=slate] at (8.3,-1.1) {cumulative $83$};
  \draw[->,gold!80!black,dashed,line width=0.8pt] (g1.north) to[bend left=22] node[above,font=\scriptsize,text=gold!70!black]{three moves reach the top with coefficient $-d_{ijk}$: the $E_6$ cubic} (g3.north);
\end{tikzpicture}
\caption{The cubic ladder. First-order moves combine in pairs to histories, and histories
combine with one more move to reach the two-dimensional top. The only way three moves
reach the top is through a triad of the cubic.}
\label{fig:ladder}
\end{figure}

\begin{physical}[What it means physically: cubic time]
Read degree as ``order in time''. A first-order move is a state update; two moves compose
into a second-order \emph{history}; and three reach the top doublet, whose coefficient is
the cubic volume $N$ spanned by the three moves. In this reading the $E_6$ cubic --- in
GUT language the Yukawa coupling --- is the rule by which three independent updates
produce a net displacement of the clock. The algebra is exact; the reading is a bridge,
and it lacks what would make it physics: a Hamiltonian, a state and an energy scale.
\tierB
\end{physical}

\begin{keyidea}
One qutrit's nine histories build $E_8$ as $80+84+84$. The four Hesse clocks are four
orthogonal $A_2$'s whose glue is the tetracode; a clock tick grades $E_8$ as
$1+56+134+56+1$, an oriented clock as $2+27+54+82+54+27+2$; and $E_8$'s roots fibre
six-to-one over the forty points.
\end{keyidea}
